% =============================================================================
% Artigo final — Aprendizado de Máquina 2026.1 (PPgEEC/UFRN, Prof. Marcelo Fernandes)
% Formato SBC. Para compilar, coloque nesta pasta o arquivo `sbc-template.sty`
% (distribuído no zip do modelo SBC fornecido pelo professor) e, se usar as
% figuras, mantenha os caminhos relativos ../projeto-*/output/*.png.
%   pdflatex artigo.tex && pdflatex artigo.tex
% =============================================================================
\documentclass[12pt]{article}

\usepackage{sbc-template}   % <- vem do zip do modelo SBC do professor
\usepackage{graphicx}
\usepackage{url}
\usepackage[utf8]{inputenc}
\usepackage[T1]{fontenc}
\usepackage[brazil]{babel}
\usepackage{amsmath}
\usepackage{booktabs}
\usepackage{multirow}

\sloppy

\title{A Seed Única Engana: Avaliação Multi-Seed e Testes Pareados\\ em Métodos Clássicos de Aprendizado de Máquina}

\author{Vitor Y. F. Freitas\inst{1}}

\address{Programa de Pós-Graduação em Engenharia Elétrica e de Computação (PPgEEC)\\
  Universidade Federal do Rio Grande do Norte (UFRN) -- Natal, RN -- Brasil
  \email{vitoryeso@outlook.com}
}

\begin{document}

\maketitle

\begin{abstract}
  Machine learning results are routinely reported as a single number obtained
  from a single random seed, a practice that conflates real effects with
  sampling noise. Using as an empirical testbed eight classical methods
  implemented mostly from scratch for course mini-projects --- K-means, linear
  regression (batch GD vs.\ LMS), logistic regression, decision trees, RBF
  networks, SVM, MLP and CNN --- we retrofit every experiment with a
  multi-seed protocol (mean $\pm$ sample standard deviation) and per-seed
  paired comparisons. Two findings emerge. First, the relevant lever to shrink
  variance depends on its source: for stable models (linear/logistic), the
  seed-to-seed spread is dominated by the binomial sampling error of the test
  set, and enlarging the test set shrinks it (from $\pm 2.5$ p.p.\ to $\pm
  0.3$ p.p.\ in our logistic experiment); for high-variance models (unpruned
  trees, CNNs), the spread comes from the model itself and more data does not
  help (the tree's std \emph{grew} from $\pm 0.021$ to $\pm 0.033$). Second,
  paired per-seed comparisons cancel shared sampling noise, reducing the
  standard deviation of the difference by $\approx 7\times$ and cleanly
  separating true ties (Gini vs.\ entropy, $t=0.11$; RGB vs.\ grayscale,
  $t=1.3$) from real but masked effects (2D vs.\ 1D CNN, $t\approx 5$;
  K-means centers in RBF networks, $t\approx 16$--$20$). We distill these
  observations into concrete reporting recommendations for ML coursework and
  research.
\end{abstract}

\begin{resumo}
  Resultados de aprendizado de máquina costumam ser reportados como um único
  número obtido com uma única seed aleatória, prática que confunde efeitos
  reais com ruído amostral. Usando como bancada empírica oito métodos
  clássicos implementados majoritariamente do zero em miniprojetos da
  disciplina --- K-means, regressão linear (GD em lote vs.\ LMS), regressão
  logística, árvore de decisão, redes RBF, SVM, MLP e CNN ---, retrofitamos
  todos os experimentos com um protocolo multi-seed (média $\pm$
  desvio-padrão amostral) e comparações pareadas por seed. Duas descobertas
  emergem. Primeira: a alavanca correta para encolher a variância depende da
  sua origem --- em modelos estáveis (linear/logística) a dispersão entre
  seeds é dominada pelo erro amostral binomial do conjunto de teste, e
  aumentar o teste a reduz (de $\pm 2{,}5$ p.p.\ para $\pm 0{,}3$ p.p.\ na
  logística); em modelos de alta variância (árvore não podada, CNN) a
  dispersão vem do próprio modelo e mais dados não ajudam (o desvio da árvore
  \emph{cresceu} de $\pm 0{,}021$ para $\pm 0{,}033$). Segunda: a comparação
  pareada por seed cancela o ruído amostral compartilhado, reduz o
  desvio-padrão da diferença em $\approx 7\times$ e separa com clareza
  empates reais (Gini vs.\ entropia, $t=0{,}11$; RGB vs.\ cinza, $t=1{,}3$)
  de efeitos reais mascarados (CNN 2D vs.\ 1D, $t\approx 5$; centros K-means
  em redes RBF, $t\approx 16$--$20$). Destilamos essas observações em
  recomendações concretas de como reportar resultados de AM.
\end{resumo}

% =============================================================================
\section{Introdução}\label{sec:intro}

Quase todo curso e boa parte da literatura aplicada de aprendizado de máquina
(AM) reporta desempenho como um número pontual: ``a regressão logística
atingiu 96{,}7\% de acurácia'', ``Gini (96\%) superou entropia (94\%)''.
Esses números, porém, são \emph{realizações de variáveis aleatórias}: dependem
da seed que gerou os dados sintéticos, do embaralhamento do split
treino/teste, da inicialização dos pesos e da ordem de apresentação dos
exemplos. Uma única realização não distingue um efeito real de um golpe de
sorte --- e, como mostraremos com experimentos deste próprio trabalho, ambos
os exemplos acima eram enganosos: o ``96{,}7\%'' era uma realização otimista
de um teste com apenas 60 amostras (o valor multi-seed é $96{,}3\% \pm
0{,}3\%$), e a ``vantagem de 2 pontos do Gini'' era puro ruído de amostragem
(o efeito pareado medido é $-0{,}0002 \pm 0{,}008$, $t=0{,}11$ --- zero).

Este artigo usa como bancada empírica os oito miniprojetos da disciplina de
Aprendizado de Máquina (PPgEEC/UFRN, 2026.1), que implementam
majoritariamente do zero uma trilha completa de métodos clássicos: K-means,
regressão linear otimizada por gradiente em lote (GD) e por LMS/Widrow-Hoff,
regressão logística, árvore de decisão CART, redes de funções de base radial
(RBF), SVM com kernel RBF, perceptron multicamadas (MLP) e redes
convolucionais (CNN). Originalmente, todos os projetos rodavam com uma seed
fixa (42) e reportavam um único número. Nossa contribuição é
\emph{retrofitar} toda a bancada com um protocolo estatístico uniforme ---
repetição multi-seed com média $\pm$ desvio-padrão amostral e, nas
comparações entre variantes de um mesmo método, teste pareado por seed --- e
analisar o que muda nas conclusões.

As perguntas de pesquisa são:

\begin{enumerate}
  \item \textbf{RQ1} --- Quais conclusões tiradas com seed única sobrevivem à
        validação multi-seed, e quais se revelam ruído?
  \item \textbf{RQ2} --- De onde vem a variância entre seeds: do erro amostral
        da \emph{medida} (teste finito) ou da instabilidade do próprio
        \emph{modelo}? A resposta muda a alavanca de mitigação correta?
  \item \textbf{RQ3} --- Quanto o pareamento por seed aumenta a sensibilidade
        das comparações entre métodos, e que decisões ele muda na prática?
\end{enumerate}

As respostas resumidas: (RQ1) três conclusões de seed única caíram e duas se
confirmaram; (RQ2) a origem da variância separa os modelos em dois regimes
com receitas opostas --- para modelos estáveis, aumentar o conjunto de teste
(de graça, quando os dados são sintéticos) encolhe o desvio; para modelos de
alta variância, só regularizar o modelo ajuda; (RQ3) o pareamento reduziu o
desvio da diferença em $\approx 7\times$ e revelou tanto empates verdadeiros
quanto um sinal real que a comparação não pareada mascarava.

O restante do artigo se organiza assim: a Seção~\ref{sec:sota} revisa a
literatura de testes estatísticos e reprodutibilidade em AM; a
Seção~\ref{sec:met} descreve a bancada (métodos, dados e protocolo); a
Seção~\ref{sec:res} apresenta os resultados; a Seção~\ref{sec:ana} discute as
duas descobertas e suas implicações; a Seção~\ref{sec:conc} conclui.

% =============================================================================
\section{Estado da Arte}\label{sec:sota}

A preocupação com o rigor estatístico na comparação de algoritmos de AM não é
nova. \cite{dietterich1998} formalizou o problema no contexto de
classificação supervisionada, mostrou que testes ingênuos (como o teste $t$
sobre repetições com dados sobrepostos) têm erro Tipo~I inflado e propôs
alternativas hoje canônicas: o teste de McNemar para uma única partição e o
teste $t$ pareado \emph{5x2cv} para comparações com re-treino. Na mesma
linha, \cite{salzberg1997} catalogou armadilhas recorrentes --- comparações
múltiplas sem correção, reuso do conjunto de teste, ausência de
replicações --- e \cite{nadeau2003} mostraram que a variância do erro de
generalização estimado por reamostragem é sistematicamente subestimada pelos
estimadores usuais, propondo correções que hoje levam seu nome.
\cite{demsar2006} estendeu a discussão para comparações de múltiplos
classificadores sobre múltiplos conjuntos de dados, recomendando testes não
paramétricos (Wilcoxon, Friedman com post-hoc de Nemenyi) no lugar de médias
de acurácia --- trabalho que se tornou referência obrigatória na área.

Uma segunda onda de interesse veio com o aprendizado profundo, no qual o
custo de cada execução tornou tentador (e comum) reportar corridas únicas.
\cite{henderson2018} mostraram que, em aprendizado por reforço profundo, a
simples mudança de seed produz curvas de recompensa tão diferentes que
subconjuntos de seeds do \emph{mesmo} algoritmo parecem algoritmos distintos
--- e que boa parte dos resultados publicados não sobrevive a re-execução.
\cite{reimers2017} documentaram o mesmo fenômeno em PLN: a diferença entre
duas arquiteturas de rede para etiquetagem de sequências é frequentemente
menor que a dispersão entre seeds de uma delas, e recomendam reportar
\emph{distribuições} de escores em vez de valores únicos.
\cite{bouthillier2021} deram o tratamento mais sistemático ao problema:
decompuseram a variância de benchmarks de AM em suas fontes (splits de
dados, inicialização, ordem dos lotes, ajuste de hiperparâmetros), mostraram
que a variância induzida pelos dados domina a da inicialização e que
randomizar o máximo possível de fontes --- em vez de fixá-las --- produz
estimativas menos enviesadas da diferença entre algoritmos. No plano
institucional, \cite{pineau2021} descrevem o programa de reprodutibilidade
do NeurIPS, que consagrou checklists exigindo múltiplas execuções e medidas
de dispersão.

Nosso trabalho não propõe um teste novo: ele transporta esse ferramental ---
repetição multi-seed \cite{reimers2017,bouthillier2021} e comparação pareada
\cite{dietterich1998,demsar2006} --- para uma bancada didática de métodos
clássicos implementados do zero, e extrai dela duas lições que a literatura
trata separadamente: a distinção entre variância da medida e variância do
modelo \cite{nadeau2003,bouthillier2021} e o ganho de sensibilidade do
pareamento por seed. A novidade está na demonstração controlada, de ponta a
ponta e com os mesmos dados e código, de como cada conclusão de seed única se
transforma (ou não) sob o protocolo estatístico.

% =============================================================================
\section{Metodologia}\label{sec:met}

\subsection{Bancada: os oito métodos}\label{sec:met-metodos}

A bancada empírica são os miniprojetos da disciplina (numerados P1--P10, sem
P6), implementados em Rust, Python/numpy ou PyTorch. Salvo indicação, os
modelos foram implementados \emph{do zero}, sem bibliotecas de AM:

\begin{itemize}
  \item \textbf{P1 --- K-means} (Rust, do zero): algoritmo de Lloyd com
        inicialização K-means++ sobre 8.000 embeddings CLIP (512-d) de uma
        coleção pessoal de imagens; escolha de $K$ por cotovelo e silhueta.
  \item \textbf{P2 --- Classificação de taxonomia descoberta} (Python):
        \emph{linear probes} sobre os rótulos derivados do clustering de P1,
        usados aqui como estudo da circularidade de avaliação
        (Seção~\ref{sec:res-outros}).
  \item \textbf{P3 --- Regressão linear: GD vs.\ LMS} (Rust, do zero):
        mesma regressão linear otimizada por gradiente em lote e pela regra
        LMS/Widrow-Hoff \cite{widrow1960} sobre dados sintéticos com pesos
        verdadeiros conhecidos ($y = 3x_1 - 2x_2 + 1 + \varepsilon$,
        $\sigma^2 = 0{,}25$).
  \item \textbf{P4 --- Regressão logística} (Rust, do zero): sigmoid
        numericamente estável, perda BCE e gradiente analítico, sobre duas
        gaussianas 2D com limite de Bayes conhecido
        ($\Phi(-\sqrt{13}/2) \approx 3{,}57\%$ de erro).
  \item \textbf{P5 --- Árvore de decisão CART} (Python, do zero): impureza,
        busca gulosa de split, recursão e predição manuais \cite{breiman1984};
        comparação dos critérios Gini vs.\ entropia.
  \item \textbf{P7 --- Rede RBF} (numpy, do zero): camada de gaussianas +
        saída linear por mínimos quadrados \cite{broomhead1988}; ablação da
        estratégia de escolha de centros (aleatório, subconjunto, K-means).
  \item \textbf{P8 --- SVM com kernel RBF} (scikit-learn): varredura
        $C \times \gamma$ \cite{cortes1995}; único modelo não implementado à
        mão (o foco era o papel dos vetores de suporte).
  \item \textbf{P9 --- MLP} (numpy, backpropagation manual
        \cite{rumelhart1986}): MSE vs.\ entropia cruzada e ablação de
        arquitetura (0 a 128 neurônios ocultos).
  \item \textbf{P10 --- CNN} (PyTorch \cite{lecun1998}): CNN 2D vs.\ 1D,
        entrada RGB vs.\ tons de cinza, ablação de número de filtros e
        tamanho de kernel, e visualização de \emph{feature maps}.
\end{itemize}

\subsection{Dados}\label{sec:met-dados}

Três famílias de dados são usadas:

\begin{itemize}
  \item \textbf{Sintéticos com verdade conhecida} (P3, P4, P5): gaussianas 2D
        e \texttt{make\_classification}. A vantagem metodológica é dupla:
        (i) o ótimo é conhecido (piso de ruído $\sigma^2$ em P3; limite de
        Bayes em P4), permitindo aferir o quão perto o modelo chega do
        possível; (ii) dados novos são \emph{de graça} --- cada seed gera uma
        realização inteiramente nova do problema, e o tamanho do conjunto de
        teste é um parâmetro livre.
  \item \textbf{Features CLIP reais} (P1, P2, P7, P8, P9): embeddings de
        512 dimensões \cite{radford2021} extraídos de $\sim$8.000 imagens de
        uma coleção pessoal, com três tarefas binárias de dificuldade
        crescente: \emph{macro\_vs\_rest} (balanceada, quase separável),
        \emph{has\_people} e \emph{screenshot\_vs\_photo} (desbalanceadas,
        base 84\%). P7--P9 usam o mesmo split, permitindo a meta-comparação
        da Seção~\ref{sec:res-outros}.
  \item \textbf{Imagens brutas} (P10): 450 imagens ($64\times 64$) da mesma
        coleção, na tarefa real foto vs.\ captura de tela.
\end{itemize}

\subsection{Protocolo multi-seed}\label{sec:met-multiseed}

Cada experimento é repetido com $N$ seeds ($N \in \{3, 5, 10, 30\}$ conforme
o custo computacional do projeto; a lista de seeds é fixa e publicada nos
relatórios). A seed controla todas as fontes de aleatoriedade do pipeline:
geração dos dados sintéticos (quando aplicável), embaralhamento do split
treino/teste, inicialização de pesos/centróides e ordem de apresentação.
Reportamos média e desvio-padrão \emph{amostral},

\begin{equation}
  \bar{m} = \frac{1}{N}\sum_{s=1}^{N} m_s, \qquad
  \hat{\sigma} = \sqrt{\frac{1}{N-1}\sum_{s=1}^{N}(m_s - \bar{m})^2},
  \label{eq:multiseed}
\end{equation}

onde $m_s$ é a métrica (acurácia, F1, MSE\ldots) obtida com a seed $s$ e o
denominador $N-1$ (\texttt{ddof=1}) corrige o viés do estimador de variância
para amostras pequenas. O número reportado como \emph{oficial} de cada
projeto é sempre $\bar{m} \pm \hat{\sigma}$; valores de seed única são
mantidos nos relatórios apenas como ilustração, com essa ressalva explícita.

\subsection{Teste pareado por seed}\label{sec:met-pareado}

Para comparar duas variantes $A$ e $B$ de um mesmo método (Gini vs.\
entropia; centros K-means vs.\ aleatórios; CNN 2D vs.\ 1D), o protocolo não
pareado --- comparar $\bar{m}_A \pm \hat{\sigma}_A$ com $\bar{m}_B \pm
\hat{\sigma}_B$, cada um sobre sorteios independentes --- desperdiça
informação: grande parte de $\hat{\sigma}$ é ruído de amostragem
\emph{compartilhado} (a mesma seed produz um split fácil ou difícil para
ambas as variantes). No protocolo pareado, cada seed $s$ gera \emph{uma
única} realização dos dados (mesmo dataset, mesmo split) sobre a qual ambas
as variantes são treinadas e avaliadas; a estatística de interesse é a
diferença por seed $d_s = m_{A,s} - m_{B,s}$, resumida por

\begin{equation}
  \bar{d} \pm \hat{\sigma}_d, \qquad
  t = \frac{\bar{d}}{\hat{\sigma}_d / \sqrt{N}},
  \label{eq:paired}
\end{equation}

onde $\bar{d}$ é a diferença média, $\hat{\sigma}_d$ o desvio-padrão amostral
das diferenças e $t$ a estatística do teste $t$ pareado (a diferença média em
unidades do seu erro-padrão). Como o ruído comum a $A$ e $B$ se cancela em
$d_s$, $\hat{\sigma}_d$ pode ser muito menor que $\hat{\sigma}_A$ ou
$\hat{\sigma}_B$ --- na prática observamos redução de $\approx 7\times$
(Seção~\ref{sec:res-p5}). Com $N=30$ seeds, $|t| \gtrsim 2$ indica efeito
distinguível de zero ao nível usual de 5\%; valores de $|t| \ll 1$ indicam
efeito \emph{medido como zero}, o que é uma afirmação mais forte do que
``não conseguimos distinguir'' \cite{dietterich1998}.

\subsection{Diagnóstico: variância da medida vs.\ variância do modelo}
\label{sec:met-var}

Por fim, para responder à RQ2, usamos um diagnóstico simples. Se o modelo é
(quase) determinístico dado o problema --- caso dos modelos lineares
convexos ---, a dispersão entre seeds da acurácia de teste deve coincidir com
o erro-padrão \emph{binomial} de medir uma proporção $p$ com $n$ amostras de
teste:

\begin{equation}
  \mathrm{EP}_{\text{binomial}} = \sqrt{\frac{p(1-p)}{n_{\text{teste}}}},
  \label{eq:binomial}
\end{equation}

que depende apenas do tamanho do teste. Nesse regime, aumentar
$n_{\text{teste}}$ encolhe o desvio como $1/\sqrt{n}$, e rodar mais seeds
apenas re-mede a mesma incerteza. Se, ao contrário, a dispersão observada
excede (ou não responde a) esse previsor, a variância vem da instabilidade
do próprio modelo --- pequenas mudanças no treino mudam a função aprendida
--- e a alavanca passa a ser regularização/poda, não mais dados de teste. O
experimento crítico é variar $n_{\text{teste}}$ e observar se
$\hat{\sigma}$ responde.

% =============================================================================
\section{Resultados}\label{sec:res}

\subsection{P4 --- Regressão logística: a seed sortuda e o teste pequeno}
\label{sec:res-p4}

Com a seed original (42) e um teste de apenas 60 amostras, a logística
reportava 96{,}67\% de acurácia de teste (erro 3{,}33\%, nominalmente
\emph{abaixo} do limite de Bayes de 3{,}57\% --- um alerta em si). O
protocolo multi-seed (10 seeds, dataset de 10.000 amostras por seed, teste
$\sim$2.000) produz o resultado oficial da Tabela~\ref{tab:p4}.

\begin{table}[htb]
\centering
\caption{P4 --- Regressão logística (do zero, Rust): multi-seed, 10 seeds,
teste $\sim$2.000 amostras por seed. Fonte:
\texttt{projeto-4/output/report\_multiseed.txt}.}
\label{tab:p4}
\begin{tabular}{lcc}
\toprule
Métrica (teste) & Média & Desvio ($\hat{\sigma}$, ddof=1) \\
\midrule
Acurácia   & 96{,}29\% & 0{,}27\% \\
Precisão   & 0{,}9609  & 0{,}0057 \\
Recall     & 0{,}9652  & 0{,}0044 \\
F1         & 0{,}9630  & 0{,}0030 \\
BCE        & 0{,}0977  & 0{,}0054 \\
\midrule
Erro de teste & \multicolumn{2}{c}{$3{,}71\% \pm 0{,}27\%$} \\
Limite de Bayes (teórico) & \multicolumn{2}{c}{$3{,}57\%$
  \ [$\Phi(-\sqrt{13}/2)$]} \\
\bottomrule
\end{tabular}
\end{table}

Três leituras. (i)~O número oficial é $96{,}3\% \pm 0{,}3\%$; o ``96{,}7\%''
da seed única era uma realização otimista dentro da variância esperada ---
entre as 10 seeds do protocolo original com teste de 60 pontos, uma chegou a
100\%. (ii)~A Eq.~\eqref{eq:binomial} explica quantitativamente o ganho: com
$n=60$, $\mathrm{EP} \approx \sqrt{0{,}96 \cdot 0{,}04/60} \approx 2{,}5$
p.p.; com $n=2.000$, $\approx 0{,}4$ p.p.\ --- e o desvio observado entre
seeds ($0{,}27$ p.p.) é compatível com esse previsor binomial, confirmando
que a logística é um modelo estável: a variância era da \emph{medida}, não do
modelo. Como os dados são sintéticos, o teste maior é gratuito. (iii)~O erro
oficial ($3{,}71\% \pm 0{,}27\%$) fica a $\sim$0,5 desvio do limite de Bayes:
o gargalo é o ruído irredutível do problema, e o modelo está essencialmente
resolvido.

\subsection{P3 --- GD vs.\ LMS: um efeito real que sobrevive}
\label{sec:res-p3}

Nem toda conclusão de seed única cai. Na regressão linear, o LMS (60.000
passos por execução) converge ao piso de ruído e o GD em lote com
\emph{learning rate} conservador (escolha didática deliberada) fica no meio
do caminho. O multi-seed (10 seeds, $N=10.000$) confirma e quantifica:

\begin{table}[htb]
\centering
\caption{P3 --- Regressão linear: GD vs.\ LMS, 10 seeds, $N=10.000$.
Verdadeiros: $w = [3, -2]$, $b = 1$; piso de ruído $\sigma^2 = 0{,}25$.
Fonte: \texttt{projeto-3/RELATORIO.md}.}
\label{tab:p3}
\begin{tabular}{lccc}
\toprule
Métrica & GD & LMS & Referência \\
\midrule
MSE final & $0{,}5018 \pm 0{,}0118$ & $0{,}2517 \pm 0{,}0039$ & 0{,}25 \\
$w_0$     & $2{,}60 \pm 0{,}013$    & $2{,}996 \pm 0{,}011$   & 3{,}0 \\
$w_1$     & $-1{,}73 \pm 0{,}012$   & $-2{,}001 \pm 0{,}006$  & $-2{,}0$ \\
$b$       & $0{,}86 \pm 0{,}014$    & $0{,}991 \pm 0{,}016$   & 1{,}0 \\
\bottomrule
\end{tabular}
\end{table}

O gap de MSE entre GD e LMS ($\approx 0{,}25$) é $\approx 20\times$ maior que
o desvio de qualquer um dos dois --- as bandas nem se aproximam. Este é o
padrão de um \emph{efeito real}: robusto, com margem folgada sobre o ruído.
Note também o comportamento da variância: ao passar de $N=300$ (seed única)
para $N=10.000$, o desvio encolheu $\approx 4\times$ --- de novo o regime de
modelo estável da Seção~\ref{sec:met-var}, em que mais dados firmam o número.

\subsection{P5 --- Gini vs.\ entropia: o empate que parecia vitória}
\label{sec:res-p5}

Este é o caso central do artigo. Com seed única (42, teste de 100 amostras),
a árvore CART manual dava a impressão de uma vitória clara do Gini
(Tabela~\ref{tab:p5-single}). O multi-seed não pareado já derruba a
conclusão; o teste pareado a substitui por uma afirmação mais forte.

\begin{table}[htb]
\centering
\caption{P5 --- Árvore de decisão, seed única (42): a ``vitória'' do Gini.
Fonte: \texttt{projeto-5/output/report.txt}.}
\label{tab:p5-single}
\begin{tabular}{lcc}
\toprule
Métrica (seed=42) & Gini & Entropia \\
\midrule
Acurácia (treino) & 1{,}0000 & 1{,}0000 \\
Acurácia (teste)  & \textbf{0{,}9600} & 0{,}9400 \\
F1                & 0{,}9608 & 0{,}9412 \\
Profundidade      & 9        & 13 \\
Folhas com 1 amostra & 10    & 8 \\
\bottomrule
\end{tabular}
\end{table}

\begin{table}[htb]
\centering
\caption{P5 --- O mesmo experimento sob os três protocolos. Não pareado:
10 seeds, $N=2.000$. Pareado: 30 seeds, mesmo split por seed, teste $N=400$.
Fonte: \texttt{projeto-5/RELATORIO.md}.}
\label{tab:p5-proto}
\begin{tabular}{llll}
\toprule
Protocolo & Gini & Entropia & Conclusão \\
\midrule
Seed única (42)     & 0{,}960 & 0{,}940 & ``Gini vence por 2 p.p.'' \\
Multi-seed não pareado & $0{,}943 \pm 0{,}033$ & $0{,}945 \pm 0{,}033$ &
  indistinguíveis (ranking inverte) \\
Pareado (30 seeds)  & \multicolumn{2}{l}{$d = -0{,}0002 \pm 0{,}008$
  \ ($t = 0{,}11$)} & \textbf{empate real: efeito $\approx$ 0} \\
\bottomrule
\end{tabular}
\end{table}

Três observações:

\begin{itemize}
  \item \textbf{O pareamento encolhe o desvio $\approx 7\times$}: o
        desvio-padrão da \emph{diferença} pareada ($\pm 0{,}008$) é cerca de
        sete vezes menor que o desvio não pareado ($\pm 0{,}055$ no mesmo
        protocolo de 30 seeds), porque o ruído de amostragem que afeta os
        dois critérios igualmente se cancela na diferença por seed. A
        Figura~\ref{fig:paired} (\texttt{projeto-5/output/paired\_viz.png})
        visualiza o contraste: à esquerda, as distribuições não pareadas se
        sobrepõem largamente; à direita, a distribuição das diferenças por
        seed se concentra em torno de zero.
  \item \textbf{$t = 0{,}11$ não é ``empate por falta de dados''}: a
        diferença média está a 0{,}1 erro-padrão de zero. Com a
        sensibilidade que o pareamento dá, um efeito de meio ponto percentual
        seria detectado; o efeito medido é zero. O único efeito real e
        consistente entre os critérios é \emph{estrutural}: a árvore da
        entropia cresce mais profunda (17{,}8 vs.\ 15{,}5 níveis sem poda),
        porque a maior curvatura da entropia perto de $p = 0{,}5$ reordena os
        splits gulosos.
  \item \textbf{Com poda, o empate persiste} ($t = 0{,}57$ com
        \texttt{max\_depth}=5) e a acurácia até melhora (0{,}947 vs.\
        0{,}936 sem poda) --- evidência adicional de que a árvore irrestrita
        (100\% no treino, folhas de 1 amostra) estava memorizando.
\end{itemize}

\begin{figure}[htb]
\centering
% Caminho relativo à pasta artigo/; ajuste se compilar de outro diretório.
\includegraphics[width=0.95\textwidth]{../projeto-5/output/paired_viz.png}
\caption{Comparação não pareada (esq.) vs.\ pareada (dir.) para Gini vs.\
entropia (P5, 30 seeds). No protocolo não pareado as distribuições de
acurácia se sobrepõem e a comparação é inconclusiva; no pareado, a
distribuição das diferenças por seed concentra-se em zero com desvio
$\approx 7\times$ menor, permitindo afirmar o empate real ($t=0{,}11$).}
\label{fig:paired}
\end{figure}

Igualmente importante é o comportamento da variância com o tamanho dos
dados: ao aumentar o dataset de $N=500$ para $N=2.000$, o desvio entre seeds
\emph{cresceu} de $\pm 0{,}021$ para $\pm 0{,}033$ --- o oposto do observado
em P3/P4. A árvore não podada é um modelo de alta variância por construção
(cresce até folhas puras), então a dispersão entre seeds reflete a
instabilidade estrutural do modelo, não o erro amostral da medida; mais
dados não firmam o número, só a poda o faria (Seção~\ref{sec:ana}).

\subsection{P7 --- RBF: vitória robusta e empate real na mesma ablação}
\label{sec:res-p7}

A rede RBF manual (features CLIP, tarefa \emph{macro\_vs\_rest}) compara três
estratégias de posicionamento dos centros. Na seed única, o ranking aparente
era K-means $>$ aleatório $>$ subconjunto em todo $k$
(Tabela~\ref{tab:p7}).

\begin{table}[htb]
\centering
\caption{P7 --- Ablação de centros da RBF (F1, \emph{macro\_vs\_rest},
seed=42) e veredito do teste pareado multi-seed. Fontes:
\texttt{projeto-7/output/report.txt} e \texttt{projeto-7/slides.html}.}
\label{tab:p7}
\begin{tabular}{lcccl}
\toprule
Estratégia & $k=10$ & $k=30$ & $k=60$ & Veredito pareado \\
\midrule
K-means     & 0{,}913 & 0{,}926 & 0{,}931 & vence, robusto
  ($t \approx 16$--$20$) \\
Aleatório   & 0{,}829 & 0{,}897 & 0{,}920 &
  \multirow{2}{*}{aleatório vs.\ subset: empate ($t = 0{,}36$)} \\
Subconjunto & 0{,}723 & 0{,}871 & 0{,}915 & \\
\bottomrule
\end{tabular}
\end{table}

O teste pareado separa as duas comparações que a seed única tratava por
igual: K-means vs.\ demais é um efeito enorme e robusto ($t \approx
16$--$20$ em todos os $k$: posicionar os centros onde os dados se concentram
importa de verdade), enquanto aleatório vs.\ subconjunto --- que na seed
única parecia uma diferença de até 10 pontos em $k=10$ --- é empate real
($t = 0{,}36$): zero medido, não falta de dado. Os números oficiais da RBF
(multi-seed, 3 seeds, $k=30$, K-means) são F1 $0{,}921 \pm 0{,}005$
(\emph{macro}), $0{,}682 \pm 0{,}010$ (\emph{has\_people}) e $0{,}668 \pm
0{,}030$ (\emph{screenshot}).

\subsection{P10 --- CNN: o pareamento revela um sinal mascarado}
\label{sec:res-p10}

O caso mais instrutivo para a RQ3. Na tarefa real foto vs.\ captura de tela
(450 imagens, $64 \times 64$), a seed única sugeria o ranking RGB $>$ cinza
$>$ 1D (Tabela~\ref{tab:p10}). Como a CNN é um modelo de alta variância
(inicialização, ordem dos lotes, poucos dados), a comparação não pareada
entre seeds tem desvios grandes --- grandes o suficiente para enterrar
\emph{todas} as diferenças como ``ruído''. O teste pareado (30 seeds, mesmos
dados e split por seed) decompõe o ranking em um efeito real e um espúrio:

\begin{table}[htb]
\centering
\caption{P10 --- CNN foto vs.\ screenshot: F1 da seed única (42) e veredito
do teste pareado (30 seeds). Fontes: \texttt{projeto-10/output/report.txt} e
\texttt{projeto-10/slides.html}.}
\label{tab:p10}
\begin{tabular}{lccl}
\toprule
Comparação & \multicolumn{2}{c}{F1 (seed=42)} & Veredito pareado (30 seeds)\\
\midrule
2D-RGB vs.\ 2D-cinza & 0{,}917 & 0{,}891 &
  empate real ($t = 1{,}3$): cor não agrega \\
2D vs.\ 1D (perfil)  & 0{,}917 & 0{,}872 &
  \textbf{efeito real} ($t \approx 5$): estrutura espacial importa \\
\bottomrule
\end{tabular}
\end{table}

A leitura da seed única (``RGB $>$ cinza $>$ 1D'') misturava um sinal
verdadeiro com ruído: a vantagem do 2D sobre o 1D é real e forte ($t \approx
5$ --- colapsar a imagem num perfil de linha descarta a estrutura espacial
que a convolução 2D explora, e os \emph{feature maps} confirmam
qualitativamente: os filtros viram detectores de borda nas caixas de UI dos
screenshots), mas a vantagem do RGB sobre o cinza é empate ($t = 1{,}3$): o
canal de cor não acrescenta nada além da estrutura espacial nessa tarefa. O
ponto fino é que, \emph{sem} pareamento, a alta variância da CNN mascarava
inclusive o efeito real --- o 2D vs.\ 1D não passaria de um teste não
pareado com poucas seeds. O pareamento não serve só para desmascarar falsas
vitórias; serve também para \emph{resgatar} efeitos verdadeiros que o ruído
esconde. Na ablação restante, o multi-seed confirma o ponto ótimo em kernel
$3$ (F1 0{,}913 na seed única) com queda forte em kernel $7$ (0{,}793) ---
filtro grande demais para imagens $64 \times 64$.

\subsection{Demais projetos e meta-comparação}\label{sec:res-outros}

Os demais projetos seguem o mesmo protocolo; destacamos os resultados que
alimentam a análise da Seção~\ref{sec:ana}:

\begin{itemize}
  \item \textbf{P8 --- SVM-RBF} (3 seeds): F1 \emph{macro} $0{,}954 \pm
        0{,}004$; \emph{has\_people} $0{,}883 \pm 0{,}010$;
        \emph{screenshot} $0{,}770 \pm 0{,}007$. O efeito de $\gamma$ é
        gigantesco e robusto ($\gamma = 0{,}1$ colapsa o F1 para 0{,}09;
        efeito $\approx 200\times$ o desvio entre seeds) --- exemplo de
        conclusão de seed única que o multi-seed confirma com folga.
  \item \textbf{P9 --- MLP} (3 seeds): F1 \emph{macro} $0{,}965 \pm 0{,}001$
        (entropia cruzada). Dois efeitos pequenos que a seed única não
        resolvia ganham veredito: CE fica consistentemente acima do MSE
        ($0{,}965 \pm 0{,}001$ vs.\ $0{,}961 \pm 0{,}003$, bandas sem
        sobreposição) e a camada oculta rende $\approx 1$ p.p.\ reprodutível
        sobre a logística ($0{,}954 \pm 0{,}002$ sem camada oculta,
        $\approx 3\hat{\sigma}$) --- ganho real, porém pequeno, pois as
        features CLIP já são quase linearmente separáveis.
  \item \textbf{P2 --- circularidade quantificada} (5 seeds): classificar os
        clusters no próprio espaço que os gerou (CLIP) dá acurácia
        $0{,}954$ (macro) / $0{,}896$ (folha); num espaço independente
        (ConvNeXt), $0{,}833$ / $0{,}661$. O gap CLIP$\to$ConvNeXt vale
        \textbf{16--47$\times$} o desvio entre seeds ($\hat{\sigma}$ de
        0{,}004 a 0{,}008) --- a circularidade é um efeito real e mensurável,
        não um acaso do split.
  \item \textbf{Meta-comparação P7/P8/P9} (mesmo dataset e split, F1
        \emph{macro\_vs\_rest}, multi-seed): MLP $0{,}965 \pm 0{,}001$ $>$
        SVM $0{,}954 \pm 0{,}004$ $>$ RBF $0{,}921 \pm 0{,}005$. Os gaps entre
        modelos são maiores que os desvios, então o ranking é estatisticamente
        sustentado; nas tarefas desbalanceadas o ranking muda (SVM na frente),
        ilustrando o \emph{No Free Lunch}.
\end{itemize}

% =============================================================================
\section{Análise}\label{sec:ana}

\subsection{Descoberta 1: a alavanca certa depende da origem da variância}
\label{sec:ana-var}

Os resultados de P3/P4 e P5/P10 formam um contraste limpo entre dois regimes
(Tabela~\ref{tab:regimes}).

\begin{table}[htb]
\centering
\caption{Dois regimes de variância entre seeds e a alavanca de mitigação
correta em cada um.}
\label{tab:regimes}
\begin{tabular}{p{0.3\textwidth}p{0.3\textwidth}p{0.3\textwidth}}
\toprule
 & \textbf{Modelo estável} (linear, logística) & \textbf{Modelo de alta
variância} (árvore não podada, CNN pequena) \\
\midrule
Origem da dispersão entre seeds & Erro amostral da \emph{medida}: teste
finito, $\mathrm{EP} \approx \sqrt{p(1-p)/n}$ & Instabilidade do próprio
\emph{modelo}: a função aprendida muda com a seed \\
Evidência na bancada & P4: $\pm 2{,}5$ p.p.\ ($n{=}60$) $\to$ $\pm 0{,}3$
p.p.\ ($n{\approx}2.000$); desvio observado $\approx$ previsor binomial.
P3: desvio encolheu $4\times$ com $N$ maior & P5: desvio \emph{cresceu} de
$\pm 0{,}021$ para $\pm 0{,}033$ ao ir de $N{=}500$ para $N{=}2.000$;
P10: desvios grandes apesar do teste fixo \\
Alavanca que funciona & Aumentar o conjunto de teste (de graça em dados
sintéticos); seeds servem só para confirmar consistência & Regularizar o
modelo (poda, \texttt{max\_depth}); em P5, podar mantém o empate e
\emph{melhora} a acurácia \\
Alavanca que \emph{não} funciona & Rodar mais seeds com o mesmo teste pequeno
(re-mede a mesma incerteza) & Aumentar os dados de teste (a variância não
está na medida) \\
\bottomrule
\end{tabular}
\end{table}

A lição prática: antes de ``rodar mais seeds'', vale perguntar \emph{o que} a
dispersão entre seeds está medindo. O diagnóstico é barato --- comparar o
desvio observado com o previsor binomial da Eq.~\eqref{eq:binomial} e/ou
variar $n_{\text{teste}}$ e ver se o desvio responde. Essa distinção é a
mesma que \cite{nadeau2003} e \cite{bouthillier2021} fazem formalmente entre
os componentes de variância do estimador do erro de generalização; nossa
bancada a exibe de forma controlada, com os dois regimes convivendo no mesmo
conjunto de projetos.

\subsection{Descoberta 2: parear por seed muda vereditos}
\label{sec:ana-pareado}

O pareamento por seed é a versão ``multi-seed'' do princípio dos testes
pareados clássicos \cite{dietterich1998}: avaliar os dois tratamentos na
mesma realização e analisar as diferenças. Nossa bancada mostra os três
desfechos possíveis, todos com decisões práticas distintas:

\begin{enumerate}
  \item \textbf{Falsa vitória $\to$ empate real} (P5 Gini vs.\ entropia,
        $t=0{,}11$; P7 aleatório vs.\ subconjunto, $t=0{,}36$; P10 RGB vs.\
        cinza, $t=1{,}3$): a escolha entre as variantes é indiferente para o
        desempenho; a decisão deve se apoiar em outros critérios (custo,
        interpretabilidade --- p.\ ex., a árvore Gini é mais rasa).
  \item \textbf{Vitória confirmada e quantificada} (P7 K-means,
        $t \approx 16$--$20$; P3 LMS vs.\ GD, gap $20\times$ o desvio): o
        efeito sobrevive com margem e pode ser reportado sem ressalva.
  \item \textbf{Sinal mascarado $\to$ efeito real} (P10 2D vs.\ 1D,
        $t \approx 5$): sem pareamento, a alta variância do modelo enterraria
        um efeito verdadeiro. Este desfecho é o menos discutido em cursos: o
        rigor estatístico não serve apenas para derrubar claims, mas para
        \emph{recuperar potência} e afirmar com segurança efeitos que uma
        comparação ingênua descartaria.
\end{enumerate}

A mecânica do ganho é a covariância: $\mathrm{Var}(m_A - m_B) =
\mathrm{Var}(m_A) + \mathrm{Var}(m_B) - 2\,\mathrm{Cov}(m_A, m_B)$. Quando as
duas variantes são avaliadas na mesma realização, a covariância é alta (um
split difícil prejudica ambas) e a variância da diferença despenca --- os
$7\times$ observados em P5 equivalem a rodar $\approx 50\times$ mais seeds no
protocolo não pareado.

\subsection{Implicações: como reportar resultados de AM}
\label{sec:ana-implicacoes}

Destilamos o protocolo em cinco recomendações, todas exercitadas na bancada:

\begin{enumerate}
  \item \textbf{Nunca reportar seed única como resultado}: todo número
        oficial deve ser $\bar{m} \pm \hat{\sigma}$ (ddof=1) sobre uma lista
        fixa e publicada de seeds; valores de seed única podem permanecer
        como ilustração, rotulados como tal.
  \item \textbf{Reportar o tamanho do teste junto com a métrica}: uma
        acurácia sobre 60 amostras carrega $\pm 2{,}5$ p.p.\ de ruído
        binomial; o leitor não tem como saber sem o $n$.
  \item \textbf{Diagnosticar a origem da variância} antes de escolher a
        mitigação (Tabela~\ref{tab:regimes}).
  \item \textbf{Parear toda comparação entre variantes}: mesma realização por
        seed, estatística sobre as diferenças, reportando $\bar{d} \pm
        \hat{\sigma}_d$ e $t$ --- tanto para evitar falsas vitórias quanto
        para não perder efeitos reais.
  \item \textbf{Distinguir ``empate medido'' de ``inconclusivo''}: $|t| \ll
        1$ com pareamento sensível autoriza afirmar efeito $\approx 0$;
        intervalos largos não pareados, não.
\end{enumerate}

\subsection{Limitações}\label{sec:ana-lim}

(i)~O número de seeds varia entre projetos (3 a 30) por custo computacional;
com 3 seeds, o próprio $\hat{\sigma}$ é impreciso e os vereditos
correspondentes (P7--P9 multi-seed) devem ser lidos como indicativos ---
as comparações decisivas usaram 30 seeds. (ii)~Usamos o teste $t$ pareado
assumindo aproximadamente normalidade das diferenças; com poucas seeds,
alternativas não paramétricas (Wilcoxon \cite{demsar2006}) seriam mais
conservadoras. (iii)~Seeds re-sorteiam dados sintéticos e splits, mas nos
projetos com dados reais fixos (P7--P10) parte da variância populacional não
é capturada --- as conclusões valem para a coleção usada. (iv)~Não aplicamos
correção para comparações múltiplas; os $|t|$ decisivos aqui são extremos
($0{,}1$--$0{,}4$ ou $5$--$20$), regiões em que a correção não muda o
veredito, mas em efeitos limítrofes ($t \approx 2$) ela seria necessária.
(v)~Validação cruzada k-fold, que quantificaria a variância estrutural dos
modelos instáveis melhor que seeds isoladas, ficou como trabalho futuro.

% =============================================================================
\section{Conclusões}\label{sec:conc}

Este trabalho retrofitou uma bancada de oito métodos clássicos de AM ---
implementados majoritariamente do zero em miniprojetos da disciplina --- com
um protocolo estatístico uniforme de repetição multi-seed e testes pareados
por seed. As contribuições são:

\begin{enumerate}
  \item \textbf{Uma auditoria empírica da prática de seed única}: das
        conclusões originais, caíram a ``vitória'' do Gini sobre a entropia
        (empate real, $t=0{,}11$), a acurácia de 96{,}7\% da logística
        (realização otimista de um teste pequeno; oficial $96{,}3\% \pm
        0{,}3\%$) e a vantagem de RGB sobre cinza na CNN ($t=1{,}3$);
        sobreviveram, com margem quantificada, o gap GD vs.\ LMS
        ($20\times$ o desvio), a superioridade dos centros K-means na RBF
        ($t \approx 16$--$20$) e a vantagem da CNN 2D sobre a 1D
        ($t \approx 5$) --- esta última só detectável \emph{com} pareamento.
  \item \textbf{A distinção operacional entre variância da medida e
        variância do modelo}, com um diagnóstico barato (previsor binomial;
        resposta do desvio ao tamanho do teste) e alavancas opostas em cada
        regime --- mais dados de teste no primeiro, regularização/poda no
        segundo.
  \item \textbf{A demonstração do valor do pareamento por seed} nos três
        desfechos possíveis (empate real, vitória confirmada, sinal
        mascarado), com redução de $\approx 7\times$ no desvio da diferença.
  \item \textbf{Um protocolo de reporte} em cinco regras, aplicável
        diretamente a trabalhos de disciplina e pesquisa.
\end{enumerate}

Como trabalhos futuros: (i)~substituir/complementar o multi-seed por
validação cruzada repetida com as correções de variância de
\cite{nadeau2003}, especialmente para os modelos instáveis; (ii)~adotar
testes não paramétricos e correção de comparações múltiplas
\cite{demsar2006} nas ablações com muitas variantes; (iii)~decompor
formalmente a variância por fonte (dados, inicialização, ordem dos lotes) no
espírito de \cite{bouthillier2021}, o que a bancada já permite ao controlar
cada fonte pela seed; e (iv)~estender o protocolo pareado às comparações
\emph{entre} métodos (meta-comparação P7/P8/P9), hoje sustentada apenas por
bandas não pareadas.

% =============================================================================
% \bibliographystyle{sbc}  % <- ativar quando migrar para BibTeX (.bib)
% As referências abaixo estão embutidas para o draft compilar sem .bib;
% na versão final, migrar para referencias.bib + \bibliography{referencias}.
\begin{thebibliography}{99}

\bibitem[Bouthillier et al. 2021]{bouthillier2021}
Bouthillier, X., Delaunay, P., Bronzi, M., Trofimov, A., Nichyporuk, B.,
Szeto, J., Sepah, N., Raff, E., Madan, K., Voleti, V., Kahou, S. E.,
Michalski, V., Serdyuk, D., Arbel, T., Pal, C., Varoquaux, G., and
Vincent, P. (2021). Accounting for variance in machine learning benchmarks.
In \emph{Proceedings of Machine Learning and Systems (MLSys)}, volume 3.

\bibitem[Breiman et al. 1984]{breiman1984}
Breiman, L., Friedman, J. H., Olshen, R. A., and Stone, C. J. (1984).
\emph{Classification and Regression Trees}. Wadsworth, Belmont, CA.

\bibitem[Broomhead and Lowe 1988]{broomhead1988}
Broomhead, D. S. and Lowe, D. (1988). Multivariable functional interpolation
and adaptive networks. \emph{Complex Systems}, 2:321--355.

\bibitem[Cortes and Vapnik 1995]{cortes1995}
Cortes, C. and Vapnik, V. (1995). Support-vector networks.
\emph{Machine Learning}, 20(3):273--297.

\bibitem[Demšar 2006]{demsar2006}
Demšar, J. (2006). Statistical comparisons of classifiers over multiple data
sets. \emph{Journal of Machine Learning Research}, 7:1--30.

\bibitem[Dietterich 1998]{dietterich1998}
Dietterich, T. G. (1998). Approximate statistical tests for comparing
supervised classification learning algorithms.
\emph{Neural Computation}, 10(7):1895--1923.

\bibitem[Henderson et al. 2018]{henderson2018}
Henderson, P., Islam, R., Bachman, P., Pineau, J., Precup, D., and
Meger, D. (2018). Deep reinforcement learning that matters. In
\emph{Proceedings of the AAAI Conference on Artificial Intelligence},
volume 32.

\bibitem[LeCun et al. 1998]{lecun1998}
LeCun, Y., Bottou, L., Bengio, Y., and Haffner, P. (1998). Gradient-based
learning applied to document recognition.
\emph{Proceedings of the IEEE}, 86(11):2278--2324.

\bibitem[Nadeau and Bengio 2003]{nadeau2003}
Nadeau, C. and Bengio, Y. (2003). Inference for the generalization error.
\emph{Machine Learning}, 52(3):239--281.

\bibitem[Pineau et al. 2021]{pineau2021}
Pineau, J., Vincent-Lamarre, P., Sinha, K., Larivière, V., Beygelzimer, A.,
d'Alché-Buc, F., Fox, E., and Larochelle, H. (2021). Improving
reproducibility in machine learning research (a report from the NeurIPS 2019
reproducibility program). \emph{Journal of Machine Learning Research},
22(164):1--20.

\bibitem[Radford et al. 2021]{radford2021}
Radford, A., Kim, J. W., Hallacy, C., Ramesh, A., Goh, G., Agarwal, S.,
Sastry, G., Askell, A., Mishkin, P., Clark, J., Krueger, G., and
Sutskever, I. (2021). Learning transferable visual models from natural
language supervision. In \emph{Proceedings of the 38th International
Conference on Machine Learning (ICML)}, pages 8748--8763.

\bibitem[Reimers and Gurevych 2017]{reimers2017}
Reimers, N. and Gurevych, I. (2017). Reporting score distributions makes a
difference: Performance study of LSTM-networks for sequence tagging. In
\emph{Proceedings of the 2017 Conference on Empirical Methods in Natural
Language Processing (EMNLP)}, pages 338--348.

\bibitem[Rumelhart et al. 1986]{rumelhart1986}
Rumelhart, D. E., Hinton, G. E., and Williams, R. J. (1986). Learning
representations by back-propagating errors.
\emph{Nature}, 323:533--536.

\bibitem[Salzberg 1997]{salzberg1997}
Salzberg, S. L. (1997). On comparing classifiers: Pitfalls to avoid and a
recommended approach. \emph{Data Mining and Knowledge Discovery},
1(3):317--328.

\bibitem[Widrow and Hoff 1960]{widrow1960}
Widrow, B. and Hoff, M. E. (1960). Adaptive switching circuits. In
\emph{IRE WESCON Convention Record}, part 4, pages 96--104.

\end{thebibliography}

\end{document}
